% !TeX root = ../main.tex
\paragraph{冻结分数档}
以 90,000 条训练记录的 WOE 未取整总分为基准，在完整同分块之间选择最接近训练十分位的位置切分。
每档右端取该档训练样本的最大分数，等于切点的记录归较低一档，两端使用开口区间。
切点保存后，将同一套分档原样应用于 30,000 条测试记录。

\paragraph{计算口径}
令 $p_k$ 和 $q_k$ 分别为训练集和测试集落入第 $k$ 档的比例，评分分布的群体稳定性指标为
\[
  \mathrm{PSI}=\sum_{k=1}^{K}(q_k-p_k)\ln\frac{q_k}{p_k}.
\]
若任一档在任一子集中人数为零，则两个子集的每一档均加 $0.5$，
再分别以各自总人数加 $0.5K$ 归一化，避免零比例进入对数。
表~\ref{tab:score-psi} 中每档均有记录，本次直接使用实际占比。

\begin{table}[H]
  \centering
  \caption{训练集与测试集的固定分数档占比及 PSI}
  \label{tab:score-psi}
  \small
  \setlength{\tabcolsep}{5pt}
  \begin{tabular}{rlrrrrr}
\toprule
档 & 未取整分数范围 & 训练人数 & 训练比例 & 测试人数 & 测试比例 & PSI 分项 \\
\midrule
1 & $(-\infty, 565.795]$ & 8,999 & 10.00\% & 2,956 & 9.85\% & 0.000021 \\
2 & $(565.795, 584.817]$ & 8,992 & 9.99\% & 2,922 & 9.74\% & 0.000064 \\
3 & $(584.817, 598.843]$ & 9,065 & 10.07\% & 3,125 & 10.42\% & 0.000116 \\
4 & $(598.843, 610.102]$ & 8,996 & 10.00\% & 3,033 & 10.11\% & 0.000013 \\
5 & $(610.102, 618.406]$ & 8,964 & 9.96\% & 3,044 & 10.15\% & 0.000035 \\
6 & $(618.406, 626.591]$ & 8,985 & 9.98\% & 2,926 & 9.75\% & 0.000054 \\
7 & $(626.591, 633.237]$ & 8,984 & 9.98\% & 3,070 & 10.23\% & 0.000062 \\
8 & $(633.237, 639.830]$ & 9,204 & 10.23\% & 3,065 & 10.22\% & 0.000000 \\
9 & $(639.830, 647.650]$ & 8,657 & 9.62\% & 2,819 & 9.40\% & 0.000052 \\
10 & $(647.650, +\infty)$ & 9,154 & 10.17\% & 3,040 & 10.13\% & 0.000001 \\
\midrule
 & 合计 & 90,000 & 100.00\% & 30,000 & 100.00\% & 0.000418 \\
\bottomrule
\end{tabular}

  \smallskip

  \parbox{0.95\linewidth}{\small 切点按三位小数展示，分档使用完整精度。PSI 合计按未舍入的各档分项求和。}
\end{table}

\paragraph{结果解释}
总 PSI 为 $0.000418$。贡献最大的第 3 档，训练和测试占比分别为 $10.07\%$ 和 $10.42\%$，
按原始比例计算相差约 $0.34\%$ 。整体来看，两个随机子集在固定分数档上的分布接近。
由于样本来自同一数据集的分层随机划分，这一比较反映的是随机留出样本之间的分布差异。
检验跨期稳定性需要按评分日期获取后续客群，并沿用同一模型和固定分数档作比较。
